\begin{table}[htbp]
\centering
\small
\caption{Rendimento potencial municipal no cenário de alto insumo e sequeiro}\label{tab:h1-capacidade}
\begin{tabular}{llrrr}
\hline
Cultura & Predominância & Municípios & Mediana & P25--P75 \\
\hline
Milho & Amazônia & 76 & 3.002 & 2.862--3.264 \\
Milho & Cerrado & 61 & 3.320 & 2.797--3.572 \\
Milho & Pantanal & 4 & 6.790 & 6.321--6.933 \\
Soja & Amazônia & 76 & 2.701 & 2.557--2.852 \\
Soja & Cerrado & 61 & 3.261 & 2.817--3.390 \\
Soja & Pantanal & 4 & 3.421 & 3.386--3.457 \\
\hline
\end{tabular}
\par\smallskip\begin{minipage}{\textwidth}\footnotesize
Fonte: elaboração própria com FAO/IIASA, GAEZ v4, clima 1961--1990 e terras de 2010; agregação municipal na malha IBGE 2023. Unidade: kg de matéria seca por hectare. Cada município entra uma vez, com peso igual. P25--P75 delimita a metade central da distribuição. O cenário não representa o rendimento realizado nem garante capacidade predeterminada em relação a todo o painel.
\end{minipage}
\end{table}
